% !PRINTED 730
\weberpage{730}
%<<CONT
class number. But the class number, since there
certainly always exists at least one class, the principal class, is
different from zero, and therefore $T$ also cannot be equal to zero.
We therefore have the theorem:

1. \gesperrt{The numbers $\tau_0, \tau_1, \ldots, \tau_{\nu'-1}$ are an independent
system of normal units of the field $H_m$.}

If $\mathfrak{E}(r)$ is a normal unit, then, if we set
\weq{7}{L_\beta = \log|\mathfrak{E}(r_\beta)|}
we have, because of (3),
\weq{8}{L_\beta = -L_{\beta+\nu'}.}
Now if $\mathfrak{E}_0, \mathfrak{E}_1, \ldots, \mathfrak{E}_{\nu'-1}$ is any system of independent
% [source note: the separator after \mathfrak{E}_0 is printed as a period-shaped or damaged comma; the mathematically required list punctuation is normalized to a comma.]
normal units, we set:
\weq{9}{L_{i,k} = \log|\mathfrak{E}_i(r_k)|, \quad L_{i,k+\nu'} = -L_{i,k},}
so that the regulator of the system of the $\mathfrak{E}_i$ takes the expression:
\weq{10}{\pm\Sigma\pm L_{0,0}L_{1,1}\cdots L_{\nu'-1,\nu'-1} = L(\mathfrak{E}_0, \mathfrak{E}_1, \ldots, \mathfrak{E}_{\nu'-1}),}
and is a number different from zero, which we denote by $T_0$.

If $\mathfrak{E}(r)$ denotes any other normal unit, we can determine
the unknowns $\xi_0, \xi_1, \ldots, \xi_{\nu'-1}$ from the $\nu'$ linear equations
 that we obtain when in
\weq{11}{L_\beta = \xi_0 L_{0,\beta} + \xi_1 L_{1,\beta} + \cdots + \xi_{\nu'-1}L_{\nu'-1,\beta}}
we let the index $\beta$ run from 0 to $\nu'-1$; and because of (8) and (9)
this equation is correct for the remaining values of $\beta$ as well,
so that from it the formula valid for all $r$ can be derived:
\weq{12}{\mathfrak{E}(r) = \pm\mathfrak{E}_0(r)^{\xi_0}\mathfrak{E}_1(r)^{\xi_1}\cdots\mathfrak{E}_{\nu'-1}(r)^{\xi_{\nu'-1}}.}
If on the other hand $m_0, m_1, \ldots, m_{\nu'-1}$ denote any integral
rational numbers whatever, then all the numbers contained in the form
\[
\pm\mathfrak{E}_0(r)^{m_0}\mathfrak{E}_1(r)^{m_1}\cdots\mathfrak{E}_{\nu'-1}(r)^{m_{\nu'-1}}
\]
are normal units of the field $H_m$. Accordingly, the considerations
that we carried out in detail in \S.\,186 can be applied to the
system of equations (10), from which it follows:

2. \gesperrt{The exponents $\xi_0, \xi_1, \ldots, \xi_{\nu'-1}$ determined
from (10) or (11) are rational numbers.}
